%%%%%%%%%%%%%%%%%%%%%%%%%%%%%%%%%%%%%%%%%%%%%%%%%%%%%%%%%%%%
% 2.8 TRANSCENDENTAL NUMBER THEORY
% The Composition of Advanced Mathematics
%%%%%%%%%%%%%%%%%%%%%%%%%%%%%%%%%%%%%%%%%%%%%%%%%%%%%%%%%%%%

\section{Transcendental Number Theory}
\label{sec:transcendental-number-theory}

\prerequisite{
Elementary Number Theory from \cref{sec:elementary-number-theory},
Diophantine Approximation from \cref{sec:diophantine-approximation},
and familiarity with algebraic numbers, polynomial equations, sequences,
series, logarithms, and complex exponentials.
}

\learningobjectives{
By the end of this section, the reader should be able to:
\begin{itemize}
    \item distinguish algebraic and transcendental numbers;
    \item explain why algebraic numbers form a countable set;
    \item explain why transcendental numbers must exist;
    \item distinguish existence proofs from explicit transcendence proofs;
    \item state Liouville's theorem and explain its transcendence consequence;
    \item recognize Liouville numbers as explicit transcendental numbers;
    \item explain why Liouville's method does not detect every transcendental number;
    \item state the Hermite--Lindemann theorem;
    \item explain why \(e\) is transcendental;
    \item state the Lindemann--Weierstrass theorem in a standard form;
    \item explain why \(\pi\) is transcendental;
    \item connect the transcendence of \(\pi\) with the impossibility of
    squaring the circle;
    \item distinguish algebraic independence from ordinary transcendence;
    \item define algebraically independent numbers;
    \item explain the relationship between linear independence and
    transcendence results;
    \item state the Gelfond--Schneider theorem;
    \item apply Gelfond--Schneider to classical examples such as
    \(2^{\sqrt2}\);
    \item explain the status of numbers such as \(e+\pi\) and \(e\pi\);
    \item understand the role of exponential and logarithmic functions in
    transcendence theory;
    \item state the Lindemann--Weierstrass theorem in exponential form;
    \item explain the conceptual content of Baker's theorem on logarithms;
    \item recognize linear forms in logarithms;
    \item understand how transcendence methods can yield effective bounds
    for Diophantine equations;
    \item distinguish proven theorems from major open conjectures;
    \item state Schanuel's conjecture conceptually;
    \item recognize connections with Diophantine approximation,
    complex analysis, algebraic number theory, and Diophantine equations.
\end{itemize}
}

%%%%%%%%%%%%%%%%%%%%%%%%%%%%%%%%%%%%%%%%%%%%%%%%%%%%%%%%%%%%
\subsection{Algebraic and Transcendental Numbers}
\label{subsec:algebraic-transcendental-numbers}
%%%%%%%%%%%%%%%%%%%%%%%%%%%%%%%%%%%%%%%%%%%%%%%%%%%%%%%%%%%%

The real and complex numbers can be divided into two fundamentally different
classes.

\begin{definitionbox}{Algebraic Number}
A complex number \(\alpha\) is \emph{algebraic} if there exists a nonzero
polynomial

\[ P(x)\in\Q[x] \]

such that

\[ \boxed{P(\alpha)=0.} \]
\end{definitionbox}

Equivalently, one may require a nonzero polynomial with integer coefficients.

Examples include

\[ 2,\qquad -\frac37,\qquad \sqrt2,\qquad i,\qquad\frac{1+\sqrt5}{2}. \]

%%%%%%%%%%%%%%%%%%%%%%%%%%%%%%%%%%%%%%%%%%%%%%%%%%%%%%%%%%%%
\subsection{Transcendental Numbers}
\label{subsec:transcendental-numbers}
%%%%%%%%%%%%%%%%%%%%%%%%%%%%%%%%%%%%%%%%%%%%%%%%%%%%%%%%%%%%

\begin{definitionbox}{Transcendental Number}
A complex number is \emph{transcendental} if it is not algebraic.
\end{definitionbox}

Thus \(\alpha\) is transcendental when there is no nonzero polynomial

\[ P(x)\in\Q[x] \]

for which

\[ P(\alpha)=0. \]

Famous examples include

\[ \boxed{e\qquad\text{and}\qquad\pi.} \]

The word \emph{transcendental} reflects the fact that these numbers lie
beyond the algebraic numbers.

%%%%%%%%%%%%%%%%%%%%%%%%%%%%%%%%%%%%%%%%%%%%%%%%%%%%%%%%%%%%
\subsection{Examples and Nonexamples}
\label{subsec:transcendental-examples-nonexamples}
%%%%%%%%%%%%%%%%%%%%%%%%%%%%%%%%%%%%%%%%%%%%%%%%%%%%%%%%%%%%

Every rational number is algebraic because

\[ \frac pq \]

is a root of

\[ qx-p. \]

Similarly,

\[ \sqrt2 \]

is algebraic because it satisfies

\[ x^2-2=0. \]

The imaginary unit \(i\) is algebraic because

\[ i^2+1=0. \]

By contrast, proving that a particular number is transcendental is generally
much harder.

%%%%%%%%%%%%%%%%%%%%%%%%%%%%%%%%%%%%%%%%%%%%%%%%%%%%%%%%%%%%
\subsection{Degree of an Algebraic Number}
\label{subsec:algebraic-number-degree-transcendence}
%%%%%%%%%%%%%%%%%%%%%%%%%%%%%%%%%%%%%%%%%%%%%%%%%%%%%%%%%%%%

If \(\alpha\) is algebraic, it has a unique monic irreducible polynomial over
\(\Q\) called its minimal polynomial.

The degree of this polynomial is called the degree of \(\alpha\).

For example,

\[ \sqrt2 \]

has minimal polynomial

\[ x^2-2, \]

so

\[ [\Q(\sqrt2):\Q]=2. \]

A transcendental number has no finite algebraic degree over \(\Q\).

%%%%%%%%%%%%%%%%%%%%%%%%%%%%%%%%%%%%%%%%%%%%%%%%%%%%%%%%%%%%
\subsection{How Many Algebraic Numbers Are There?}
\label{subsec:countability-algebraic-numbers}
%%%%%%%%%%%%%%%%%%%%%%%%%%%%%%%%%%%%%%%%%%%%%%%%%%%%%%%%%%%%

The set of algebraic numbers is countable.

To see why, consider polynomials

\[ a_0+a_1x+\cdots+a_nx^n \]

with integer coefficients.

For each fixed degree and bounded coefficient size, only finitely many such
polynomials exist.

Taking the union over all degrees and coefficient bounds gives a countable
set of integer polynomials.

Each nonzero polynomial has only finitely many roots.

Therefore the set of all algebraic numbers is countable.

%%%%%%%%%%%%%%%%%%%%%%%%%%%%%%%%%%%%%%%%%%%%%%%%%%%%%%%%%%%%
\subsection{Existence of Transcendental Numbers}
\label{subsec:existence-transcendental-numbers}
%%%%%%%%%%%%%%%%%%%%%%%%%%%%%%%%%%%%%%%%%%%%%%%%%%%%%%%%%%%%

The real numbers are uncountable, while the algebraic numbers are countable.

Therefore:

\begin{theorem}
Transcendental real numbers exist.
\end{theorem}

In fact, almost every real number is transcendental in the cardinality sense.

\[
\boxed{
\R\text{ is uncountable}
\quad\text{but}\quad
\overline{\Q}\cap\R\text{ is countable}.
}
\]

Here

\[ \overline{\Q} \]

denotes the field of algebraic numbers.

The notation is read ``Q bar.''

%%%%%%%%%%%%%%%%%%%%%%%%%%%%%%%%%%%%%%%%%%%%%%%%%%%%%%%%%%%%
\subsection{Existence Versus Identification}
\label{subsec:existence-versus-identification-transcendence}
%%%%%%%%%%%%%%%%%%%%%%%%%%%%%%%%%%%%%%%%%%%%%%%%%%%%%%%%%%%%

The counting argument proves that transcendental numbers exist in enormous
abundance.

However, it does not tell us whether a specific familiar number is
transcendental.

For example, cardinality alone cannot prove that

\[ e,\qquad\pi,\qquad e^\pi \]

are transcendental.

Explicit transcendence requires additional arithmetic structure.

\begin{conceptbox}
Showing that transcendental numbers exist is relatively easy.

Showing that a particular naturally occurring number is transcendental can
be extremely difficult.
\end{conceptbox}

%%%%%%%%%%%%%%%%%%%%%%%%%%%%%%%%%%%%%%%%%%%%%%%%%%%%%%%%%%%%
\subsection{The First Explicit Transcendental Numbers}
\label{subsec:first-explicit-transcendentals}
%%%%%%%%%%%%%%%%%%%%%%%%%%%%%%%%%%%%%%%%%%%%%%%%%%%%%%%%%%%%

Liouville discovered the first systematic method for constructing explicit
transcendental numbers.

The key idea was to compare rational approximation with algebraic degree.

Algebraic numbers cannot be approximated ``too well'' by rational numbers.

Therefore a number with impossibly strong rational approximations cannot be
algebraic.

%%%%%%%%%%%%%%%%%%%%%%%%%%%%%%%%%%%%%%%%%%%%%%%%%%%%%%%%%%%%
\subsection{Liouville's Approximation Theorem}
\label{subsec:liouville-theorem-transcendence}
%%%%%%%%%%%%%%%%%%%%%%%%%%%%%%%%%%%%%%%%%%%%%%%%%%%%%%%%%%%%

The theorem was introduced in
\cref{sec:diophantine-approximation}; here we emphasize its transcendence
consequence.

\begin{theorem}[Liouville]
Let \(\alpha\) be an algebraic irrational number of degree \(d\geq2\).
Then there exists a constant \(C(\alpha)>0\) such that

\[ \boxed{\left|\alpha-\frac pq\right|>\frac{C(\alpha)}{q^d}} \]

for all rational numbers \(p/q\) with \(q>0\).
\end{theorem}

Thus algebraic numbers of fixed degree cannot possess arbitrarily strong
rational approximations.

%%%%%%%%%%%%%%%%%%%%%%%%%%%%%%%%%%%%%%%%%%%%%%%%%%%%%%%%%%%%
\subsection{Why Liouville's Bound Holds}
\label{subsec:liouville-bound-idea}
%%%%%%%%%%%%%%%%%%%%%%%%%%%%%%%%%%%%%%%%%%%%%%%%%%%%%%%%%%%%

Suppose

\[ P(x)=a_dx^d+\cdots+a_0 \]

is an integer polynomial satisfying

\[ P(\alpha)=0. \]

For a rational number

\[ \frac pq\neq\alpha, \]

the quantity

\[ q^dP\left(\frac pq\right) \]

is an integer.

Since it is nonzero,

\[ \left|q^dP\left(\frac pq\right)\right|\geq1. \]

Hence

\[ \left|P\left(\frac pq\right)\right|\geq\frac1{q^d}. \]

Near \(\alpha\), the polynomial cannot change arbitrarily quickly.

Using the Mean Value Theorem gives a constant \(M>0\) for which

\[ \left|P\left(\frac pq\right)-P(\alpha)\right|\leq M\left|\frac pq-\alpha\right|. \]

Since \(P(\alpha)=0\),

\[ \frac1{q^d}\leq M\left|\frac pq-\alpha\right|. \]

Therefore

\[ \left|\alpha-\frac pq\right|\geq\frac1{Mq^d}. \]

This is the essential mechanism behind Liouville's theorem.

%%%%%%%%%%%%%%%%%%%%%%%%%%%%%%%%%%%%%%%%%%%%%%%%%%%%%%%%%%%%
\subsection{Liouville Numbers}
\label{subsec:liouville-numbers-transcendence}
%%%%%%%%%%%%%%%%%%%%%%%%%%%%%%%%%%%%%%%%%%%%%%%%%%%%%%%%%%%%

Recall that a Liouville number is a real number \(\alpha\) such that for
every positive integer \(n\), there exist integers \(p\) and \(q>1\) with

\[ 0<\left|\alpha-\frac pq\right|<\frac1{q^n}. \]

Such approximation is stronger than the restriction permitted for any
algebraic number of fixed degree.

Therefore:

\begin{theorem}
Every Liouville number is transcendental.
\end{theorem}

%%%%%%%%%%%%%%%%%%%%%%%%%%%%%%%%%%%%%%%%%%%%%%%%%%%%%%%%%%%%
\subsection{Liouville's Constant}
\label{subsec:liouville-constant-transcendence}
%%%%%%%%%%%%%%%%%%%%%%%%%%%%%%%%%%%%%%%%%%%%%%%%%%%%%%%%%%%%

A classical example is

\[ \boxed{L=\sum_{n=1}^{\infty}10^{-n!}.} \]

Its decimal expansion begins

\[ L=0.110001000000000000000001\ldots \]

The nonzero digits occur in positions

\[ 1!,2!,3!,4!,\ldots \]

after the decimal point.

Truncating after the \(m\)-th nonzero term gives a rational number

\[ L_m=\sum_{n=1}^{m}10^{-n!}. \]

Its denominator divides

\[ 10^{m!}. \]

%%%%%%%%%%%%%%%%%%%%%%%%%%%%%%%%%%%%%%%%%%%%%%%%%%%%%%%%%%%%
\subsection{Why Liouville's Constant Is Transcendental}
\label{subsec:liouville-constant-proof-idea}
%%%%%%%%%%%%%%%%%%%%%%%%%%%%%%%%%%%%%%%%%%%%%%%%%%%%%%%%%%%%

The approximation error satisfies

\[ 0<L-L_m=\sum_{n=m+1}^{\infty}10^{-n!}. \]

The first omitted term is

\[ 10^{-(m+1)!}. \]

Since

\[ (m+1)!=(m+1)m!, \]

the error is approximately of size

\[ \frac1{(10^{m!})^{m+1}}. \]

Thus the rational approximations become stronger than any fixed power of
their denominators.

Therefore \(L\) is a Liouville number and hence transcendental.

%%%%%%%%%%%%%%%%%%%%%%%%%%%%%%%%%%%%%%%%%%%%%%%%%%%%%%%%%%%%
\subsection{Worked Example: Liouville Truncation}
\label{subsec:worked-liouville-truncation}
%%%%%%%%%%%%%%%%%%%%%%%%%%%%%%%%%%%%%%%%%%%%%%%%%%%%%%%%%%%%

\begin{workedexamplebox}{First Four Terms of \(L\)}

Compute the first four terms of

\[ L=\sum_{n=1}^{\infty}10^{-n!}. \]

\begin{tabularx}{\textwidth}{@{}>{\raggedright\arraybackslash}p{0.43\textwidth}X@{}}
Write the factorial exponents. & \(1!,2!,3!,4!=1,2,6,24\)\\
Substitute. & \(10^{-1}+10^{-2}+10^{-6}+10^{-24}\)\\
Combine the first three visible decimal terms. & \(0.110001\)\\
Add the fourth term. & \(0.110001000000000000000001\)
\end{tabularx}

\begin{answerbox}
\[ \boxed{L_4=10^{-1}+10^{-2}+10^{-6}+10^{-24}.} \]
\end{answerbox}
\end{workedexamplebox}

%%%%%%%%%%%%%%%%%%%%%%%%%%%%%%%%%%%%%%%%%%%%%%%%%%%%%%%%%%%%
\subsection{Limitations of Liouville's Method}
\label{subsec:limitations-liouville-method}
%%%%%%%%%%%%%%%%%%%%%%%%%%%%%%%%%%%%%%%%%%%%%%%%%%%%%%%%%%%%

Liouville's theorem gives an elegant transcendence criterion, but it detects
only numbers with exceptionally strong rational approximations.

Many transcendental numbers are not Liouville numbers.

In particular, transcendence does not imply

\[ \mu(\alpha)=\infty. \]

Therefore new techniques are required to prove the transcendence of familiar
constants such as \(e\) and \(\pi\).

%%%%%%%%%%%%%%%%%%%%%%%%%%%%%%%%%%%%%%%%%%%%%%%%%%%%%%%%%%%%
\subsection{The Number \(e\)}
\label{subsec:e-transcendence}
%%%%%%%%%%%%%%%%%%%%%%%%%%%%%%%%%%%%%%%%%%%%%%%%%%%%%%%%%%%%

Euler's number is

\[ e=\sum_{n=0}^{\infty}\frac1{n!}. \]

Numerically,

\[ e\approx2.718281828459045\ldots \]

The number \(e\) is irrational, but transcendence is a much stronger
statement.

\begin{theorem}[Hermite]
The number

\[ \boxed{e} \]

is transcendental.
\end{theorem}

This was the first proof that a major naturally occurring mathematical
constant is transcendental.

%%%%%%%%%%%%%%%%%%%%%%%%%%%%%%%%%%%%%%%%%%%%%%%%%%%%%%%%%%%%
\subsection{Hermite--Lindemann Theorem}
\label{subsec:hermite-lindemann}
%%%%%%%%%%%%%%%%%%%%%%%%%%%%%%%%%%%%%%%%%%%%%%%%%%%%%%%%%%%%

A stronger statement places the transcendence of \(e\) into a general
framework.

\begin{theorem}[Hermite--Lindemann]
If \(\alpha\neq0\) is algebraic, then

\[ \boxed{e^\alpha\text{ is transcendental}.} \]
\end{theorem}

For example, taking

\[ \alpha=1 \]

gives

\[ e^1=e, \]

so \(e\) is transcendental.

%%%%%%%%%%%%%%%%%%%%%%%%%%%%%%%%%%%%%%%%%%%%%%%%%%%%%%%%%%%%
\subsection{Examples from Hermite--Lindemann}
\label{subsec:hermite-lindemann-examples}
%%%%%%%%%%%%%%%%%%%%%%%%%%%%%%%%%%%%%%%%%%%%%%%%%%%%%%%%%%%%

Since

\[ 2,\qquad -3,\qquad\sqrt2,\qquad i \]

are nonzero algebraic numbers, the theorem implies that

\[
e^2,\qquad
e^{-3},\qquad
e^{\sqrt2},\qquad
e^i
\]

are transcendental.

Notice how strongly this links algebraic inputs to transcendental outputs.

%%%%%%%%%%%%%%%%%%%%%%%%%%%%%%%%%%%%%%%%%%%%%%%%%%%%%%%%%%%%
\subsection{Euler's Identity and \(\pi\)}
\label{subsec:euler-identity-pi}
%%%%%%%%%%%%%%%%%%%%%%%%%%%%%%%%%%%%%%%%%%%%%%%%%%%%%%%%%%%%

Euler's identity states

\[ \boxed{e^{i\pi}=-1.} \]

This equation creates a remarkable path toward the transcendence of \(\pi\).

Suppose, for contradiction, that \(\pi\) were algebraic.

Since \(i\) is algebraic, the product

\[ i\pi \]

would also be algebraic and nonzero.

Hermite--Lindemann would then imply

\[ e^{i\pi} \]

is transcendental.

But

\[ e^{i\pi}=-1, \]

and \(-1\) is algebraic.

Contradiction.

Therefore:

\begin{theorem}[Lindemann]
The number

\[ \boxed{\pi} \]

is transcendental.
\end{theorem}

%%%%%%%%%%%%%%%%%%%%%%%%%%%%%%%%%%%%%%%%%%%%%%%%%%%%%%%%%%%%
\subsection{Why the Transcendence of \(\pi\) Matters}
\label{subsec:why-pi-transcendence-matters}
%%%%%%%%%%%%%%%%%%%%%%%%%%%%%%%%%%%%%%%%%%%%%%%%%%%%%%%%%%%%

The transcendence of \(\pi\) settled an ancient geometric problem:
\emph{squaring the circle}.

The problem asks whether one can construct, using only an unmarked
straightedge and compass, a square having the same area as a given circle.

For a unit circle, the area is

\[ \pi. \]

A square with this area must have side length

\[ \sqrt\pi. \]

Straightedge-and-compass constructions produce only algebraic numbers of
special degree.

Since \(\pi\) is transcendental,

\[ \sqrt\pi \]

is also transcendental.

Therefore the required construction is impossible.

%%%%%%%%%%%%%%%%%%%%%%%%%%%%%%%%%%%%%%%%%%%%%%%%%%%%%%%%%%%%
\subsection{The Classical Construction Problems}
\label{subsec:classical-construction-problems-transcendence}
%%%%%%%%%%%%%%%%%%%%%%%%%%%%%%%%%%%%%%%%%%%%%%%%%%%%%%%%%%%%

Three famous Greek construction problems are:

\begin{itemize}
    \item doubling the cube;
    \item trisecting an arbitrary angle;
    \item squaring the circle.
\end{itemize}

The first two are obstructed by algebraic degree considerations.

Squaring the circle requires the stronger fact that

\[ \pi \]

is transcendental.

Thus transcendental number theory resolved a problem originating in ancient
geometry.

%%%%%%%%%%%%%%%%%%%%%%%%%%%%%%%%%%%%%%%%%%%%%%%%%%%%%%%%%%%%
\subsection{Lindemann--Weierstrass Theorem}
\label{subsec:lindemann-weierstrass}
%%%%%%%%%%%%%%%%%%%%%%%%%%%%%%%%%%%%%%%%%%%%%%%%%%%%%%%%%%%%

The Hermite--Lindemann theorem belongs to a much stronger result.

\begin{theorem}[Lindemann--Weierstrass]
If

\[ \alpha_1,\ldots,\alpha_n \]

are distinct algebraic numbers, then

\[ \boxed{e^{\alpha_1},\ldots,e^{\alpha_n}} \]

are linearly independent over the algebraic numbers.
\end{theorem}

That is, if algebraic numbers

\[ \beta_1,\ldots,\beta_n \]

satisfy

\[ \beta_1e^{\alpha_1}+\cdots+\beta_ne^{\alpha_n}=0, \]

then

\[ \beta_1=\cdots=\beta_n=0. \]

%%%%%%%%%%%%%%%%%%%%%%%%%%%%%%%%%%%%%%%%%%%%%%%%%%%%%%%%%%%%
\subsection{Linear Independence}
\label{subsec:linear-independence-transcendence}
%%%%%%%%%%%%%%%%%%%%%%%%%%%%%%%%%%%%%%%%%%%%%%%%%%%%%%%%%%%%

The phrase \emph{linearly independent over the algebraic numbers} means that
no nontrivial algebraic linear combination vanishes.

This is stronger than merely saying that each individual exponential is
transcendental.

For example,

\[ e^{\alpha_1} \]

could be transcendental while still participating in an algebraic linear
relation with other transcendental numbers.

Lindemann--Weierstrass rules out such relations for exponentials of distinct
algebraic numbers.

%%%%%%%%%%%%%%%%%%%%%%%%%%%%%%%%%%%%%%%%%%%%%%%%%%%%%%%%%%%%
\subsection{A Common Equivalent Form}
\label{subsec:lindemann-weierstrass-equivalent-form}
%%%%%%%%%%%%%%%%%%%%%%%%%%%%%%%%%%%%%%%%%%%%%%%%%%%%%%%%%%%%

Another standard formulation states that if

\[ \alpha_1,\ldots,\alpha_n \]

are algebraic numbers linearly independent over \(\Q\), then

\[ e^{\alpha_1},\ldots,e^{\alpha_n} \]

are algebraically independent in certain limited linear senses, and in
particular satisfy strong algebraic independence consequences for
appropriate monomials.

\begin{warningbox}
Linear independence and algebraic independence are different concepts.

The Lindemann--Weierstrass theorem should not be interpreted as saying that
arbitrary exponentials of algebraic numbers are algebraically independent.
\end{warningbox}

%%%%%%%%%%%%%%%%%%%%%%%%%%%%%%%%%%%%%%%%%%%%%%%%%%%%%%%%%%%%
\subsection{Algebraic Independence}
\label{subsec:algebraic-independence}
%%%%%%%%%%%%%%%%%%%%%%%%%%%%%%%%%%%%%%%%%%%%%%%%%%%%%%%%%%%%

Transcendence concerns one number at a time.

Algebraic independence concerns several numbers simultaneously.

\begin{definitionbox}{Algebraic Independence}
Numbers

\[ \alpha_1,\ldots,\alpha_n \]

are \emph{algebraically independent over \(\Q\)} if there is no nonzero
polynomial

\[ P(x_1,\ldots,x_n)\in\Q[x_1,\ldots,x_n] \]

such that

\[ \boxed{P(\alpha_1,\ldots,\alpha_n)=0.} \]
\end{definitionbox}

If such a polynomial exists, the numbers are \emph{algebraically dependent}.

%%%%%%%%%%%%%%%%%%%%%%%%%%%%%%%%%%%%%%%%%%%%%%%%%%%%%%%%%%%%
\subsection{Transcendence Degree}
\label{subsec:transcendence-degree}
%%%%%%%%%%%%%%%%%%%%%%%%%%%%%%%%%%%%%%%%%%%%%%%%%%%%%%%%%%%%

The maximum number of algebraically independent elements needed to describe
a field extension leads to the concept of transcendence degree.

The notation

\[ \operatorname{trdeg}_{\Q}K \]

is read ``the transcendence degree of \(K\) over \(\Q\).''

For example,

\[ \Q(t) \]

where \(t\) is transcendental has

\[ \operatorname{trdeg}_{\Q}\Q(t)=1. \]

This concept is developed more systematically in field theory and algebraic
geometry.

%%%%%%%%%%%%%%%%%%%%%%%%%%%%%%%%%%%%%%%%%%%%%%%%%%%%%%%%%%%%
\subsection{Transcendence Does Not Imply Independence}
\label{subsec:transcendence-not-independence}
%%%%%%%%%%%%%%%%%%%%%%%%%%%%%%%%%%%%%%%%%%%%%%%%%%%%%%%%%%%%

Suppose \(\alpha\) is transcendental.

Then

\[ \alpha\qquad\text{and}\qquad\alpha^2 \]

are both transcendental.

But they are algebraically dependent because

\[ y-x^2=0 \]

when

\[ x=\alpha,\qquad y=\alpha^2. \]

Therefore:

\[
\boxed{
\text{individual transcendence}
\not\Longrightarrow
\text{algebraic independence}.
}
\]

%%%%%%%%%%%%%%%%%%%%%%%%%%%%%%%%%%%%%%%%%%%%%%%%%%%%%%%%%%%%
\subsection{The Gelfond--Schneider Theorem}
\label{subsec:gelfond-schneider}
%%%%%%%%%%%%%%%%%%%%%%%%%%%%%%%%%%%%%%%%%%%%%%%%%%%%%%%%%%%%

Another major transcendence theorem concerns powers.

\begin{theorem}[Gelfond--Schneider]
Let \(\alpha\) and \(\beta\) be algebraic numbers such that

\[ \alpha\neq0,\qquad\alpha\neq1, \]

and suppose \(\beta\) is irrational.

Then every value of

\[ \boxed{\alpha^\beta} \]

is transcendental.
\end{theorem}

For positive real \(\alpha\), the usual real value is unambiguous.

%%%%%%%%%%%%%%%%%%%%%%%%%%%%%%%%%%%%%%%%%%%%%%%%%%%%%%%%%%%%
\subsection{The Gelfond--Schneider Constant}
\label{subsec:gelfond-schneider-constant}
%%%%%%%%%%%%%%%%%%%%%%%%%%%%%%%%%%%%%%%%%%%%%%%%%%%%%%%%%%%%

Take

\[ \alpha=2,\qquad\beta=\sqrt2. \]

Both numbers are algebraic, and \(\sqrt2\) is irrational.

Therefore

\[ \boxed{2^{\sqrt2}} \]

is transcendental.

This number is sometimes called the Gelfond--Schneider constant.

%%%%%%%%%%%%%%%%%%%%%%%%%%%%%%%%%%%%%%%%%%%%%%%%%%%%%%%%%%%%
\subsection{Worked Example: Applying Gelfond--Schneider}
\label{subsec:worked-gelfond-schneider}
%%%%%%%%%%%%%%%%%%%%%%%%%%%%%%%%%%%%%%%%%%%%%%%%%%%%%%%%%%%%

\begin{workedexamplebox}{Is \(3^{\sqrt5}\) Transcendental?}

Determine whether the Gelfond--Schneider theorem applies to

\[ 3^{\sqrt5}. \]

\begin{tabularx}{\textwidth}{@{}>{\raggedright\arraybackslash}p{0.43\textwidth}X@{}}
Check the base. & \(3\) is algebraic and \(3\neq0,1\).\\
Check the exponent. & \(\sqrt5\) is algebraic and irrational.\\
Apply Gelfond--Schneider. & \(3^{\sqrt5}\) is transcendental.
\end{tabularx}

\begin{answerbox}
\[ \boxed{3^{\sqrt5}\text{ is transcendental}.} \]
\end{answerbox}
\end{workedexamplebox}

%%%%%%%%%%%%%%%%%%%%%%%%%%%%%%%%%%%%%%%%%%%%%%%%%%%%%%%%%%%%
\subsection{Hilbert's Seventh Problem}
\label{subsec:hilbert-seventh-problem}
%%%%%%%%%%%%%%%%%%%%%%%%%%%%%%%%%%%%%%%%%%%%%%%%%%%%%%%%%%%%

The question answered by Gelfond--Schneider arose from Hilbert's seventh
problem.

The problem asked about the arithmetic nature of numbers of the form

\[ \alpha^\beta \]

when

\[
\alpha\text{ is algebraic},\qquad
\beta\text{ is algebraic irrational}.
\]

The Gelfond--Schneider theorem provides the definitive answer:
such numbers are transcendental under the theorem's hypotheses.

%%%%%%%%%%%%%%%%%%%%%%%%%%%%%%%%%%%%%%%%%%%%%%%%%%%%%%%%%%%%
\subsection{A Curious Power Identity}
\label{subsec:irrational-power-rational}
%%%%%%%%%%%%%%%%%%%%%%%%%%%%%%%%%%%%%%%%%%%%%%%%%%%%%%%%%%%%

The number

\[ \left(\sqrt2\right)^{\sqrt2} \]

is transcendental by Gelfond--Schneider.

Yet

\[
\left(
\left(\sqrt2\right)^{\sqrt2}
\right)^{\sqrt2}
=
\left(\sqrt2\right)^2
=
2.
\]

Thus a transcendental number raised to an irrational power can produce an
algebraic number.

This illustrates why arithmetic behavior under exponentiation can be
surprising.

%%%%%%%%%%%%%%%%%%%%%%%%%%%%%%%%%%%%%%%%%%%%%%%%%%%%%%%%%%%%
\subsection{The Number \(e^\pi\)}
\label{subsec:gelfond-constant}
%%%%%%%%%%%%%%%%%%%%%%%%%%%%%%%%%%%%%%%%%%%%%%%%%%%%%%%%%%%%

Euler's identity allows

\[ e^\pi \]

to be rewritten in a form suitable for Gelfond--Schneider.

Since

\[ i^{-2i}=e^\pi \]

for a suitable interpretation of the complex power, and

\[ i \]

and

\[ -2i \]

are algebraic while \(-2i\) is irrational, the theorem implies

\[ \boxed{e^\pi\text{ is transcendental}.} \]

The number \(e^\pi\) is often called Gelfond's constant.

%%%%%%%%%%%%%%%%%%%%%%%%%%%%%%%%%%%%%%%%%%%%%%%%%%%%%%%%%%%%
\subsection{What About \(\pi^e\)?}
\label{subsec:pi-to-e-status}
%%%%%%%%%%%%%%%%%%%%%%%%%%%%%%%%%%%%%%%%%%%%%%%%%%%%%%%%%%%%

The expression

\[ \pi^e \]

looks similar to

\[ e^\pi. \]

However, Gelfond--Schneider does not apply because both the base \(\pi\) and
the exponent \(e\) are transcendental.

The theorem requires an algebraic base and an algebraic irrational exponent.

This illustrates an important lesson:

\begin{conceptbox}
A theorem about transcendental numbers is only as strong as its hypotheses.
Similar-looking expressions may require completely different methods.
\end{conceptbox}

%%%%%%%%%%%%%%%%%%%%%%%%%%%%%%%%%%%%%%%%%%%%%%%%%%%%%%%%%%%%
\subsection{What Is Known About \(e+\pi\)?}
\label{subsec:e-plus-pi}
%%%%%%%%%%%%%%%%%%%%%%%%%%%%%%%%%%%%%%%%%%%%%%%%%%%%%%%%%%%%

Both

\[ e\qquad\text{and}\qquad\pi \]

are transcendental.

However, this does not automatically imply that

\[ e+\pi \]

is transcendental.

Indeed, the sum of two transcendental numbers can be algebraic.

For example, if \(\alpha\) is transcendental, then

\[ \alpha+(1-\alpha)=1. \]

Thus individual transcendence does not control sums.

%%%%%%%%%%%%%%%%%%%%%%%%%%%%%%%%%%%%%%%%%%%%%%%%%%%%%%%%%%%%
\subsection{What Is Known About \(e\pi\)?}
\label{subsec:e-times-pi}
%%%%%%%%%%%%%%%%%%%%%%%%%%%%%%%%%%%%%%%%%%%%%%%%%%%%%%%%%%%%

Likewise, the fact that \(e\) and \(\pi\) are transcendental does not by
itself determine the arithmetic nature of

\[ e\pi. \]

Products of transcendental numbers can be algebraic.

For example,

\[ \alpha\cdot\frac1\alpha=1 \]

for every nonzero transcendental \(\alpha\).

Therefore separate transcendence results cannot simply be multiplied
together.

%%%%%%%%%%%%%%%%%%%%%%%%%%%%%%%%%%%%%%%%%%%%%%%%%%%%%%%%%%%%
\subsection{A Useful Logical Principle}
\label{subsec:transcendence-logical-principle}
%%%%%%%%%%%%%%%%%%%%%%%%%%%%%%%%%%%%%%%%%%%%%%%%%%%%%%%%%%%%

Suppose \(\alpha\neq0\) is algebraic and \(\beta\) is transcendental.

Then

\[ \alpha+\beta \]

must be transcendental.

Otherwise, if \(\alpha+\beta\) were algebraic, then

\[ \beta=(\alpha+\beta)-\alpha \]

would be algebraic, a contradiction.

Similarly,

\[ \alpha\beta \]

is transcendental when \(\alpha\neq0\) is algebraic and \(\beta\) is
transcendental.

%%%%%%%%%%%%%%%%%%%%%%%%%%%%%%%%%%%%%%%%%%%%%%%%%%%%%%%%%%%%
\subsection{Worked Example: Algebraic Plus Transcendental}
\label{subsec:worked-algebraic-plus-transcendental}
%%%%%%%%%%%%%%%%%%%%%%%%%%%%%%%%%%%%%%%%%%%%%%%%%%%%%%%%%%%%

\begin{workedexamplebox}{Is \(5+\pi\) Transcendental?}

\begin{tabularx}{\textwidth}{@{}>{\raggedright\arraybackslash}p{0.43\textwidth}X@{}}
Identify \(5\). & \(5\) is algebraic.\\
Identify \(\pi\). & \(\pi\) is transcendental.\\
Suppose \(5+\pi\) were algebraic. & Then \(\pi=(5+\pi)-5\) would be algebraic.\\
Contradiction. & Therefore \(5+\pi\) is transcendental.
\end{tabularx}

\begin{answerbox}
\[ \boxed{5+\pi\text{ is transcendental}.} \]
\end{answerbox}
\end{workedexamplebox}

%%%%%%%%%%%%%%%%%%%%%%%%%%%%%%%%%%%%%%%%%%%%%%%%%%%%%%%%%%%%
\subsection{Logarithms of Algebraic Numbers}
\label{subsec:logarithms-algebraic-numbers}
%%%%%%%%%%%%%%%%%%%%%%%%%%%%%%%%%%%%%%%%%%%%%%%%%%%%%%%%%%%%

Transcendence theory also studies logarithms.

Suppose

\[ \alpha\neq0,1 \]

is algebraic and

\[ \beta=\log\alpha \]

is a nonzero logarithm of \(\alpha\).

Then

\[ e^\beta=\alpha. \]

If \(\beta\) were algebraic and nonzero, Hermite--Lindemann would imply that

\[ e^\beta \]

is transcendental.

But \(e^\beta=\alpha\) is algebraic.

Therefore every nonzero logarithm of a nonzero algebraic number other than
\(1\) is transcendental.

%%%%%%%%%%%%%%%%%%%%%%%%%%%%%%%%%%%%%%%%%%%%%%%%%%%%%%%%%%%%
\subsection{Linear Forms in Logarithms}
\label{subsec:linear-forms-logarithms}
%%%%%%%%%%%%%%%%%%%%%%%%%%%%%%%%%%%%%%%%%%%%%%%%%%%%%%%%%%%%

A central object in modern transcendence theory is an expression

\[ \Lambda=b_1\log\alpha_1+\cdots+b_n\log\alpha_n. \]

The capital Greek letter

\[ \Lambda \]

is \emph{Lambda}, pronounced ``lambda.''

Such an expression is called a \emph{linear form in logarithms}.

Typically:

\[
\alpha_1,\ldots,\alpha_n
\]

are algebraic numbers and

\[
b_1,\ldots,b_n
\]

are integers or algebraic numbers.

%%%%%%%%%%%%%%%%%%%%%%%%%%%%%%%%%%%%%%%%%%%%%%%%%%%%%%%%%%%%
\subsection{Why Linear Forms in Logarithms Matter}
\label{subsec:why-linear-forms-logarithms}
%%%%%%%%%%%%%%%%%%%%%%%%%%%%%%%%%%%%%%%%%%%%%%%%%%%%%%%%%%%%

Many exponential Diophantine equations can be rearranged into statements
that a linear form in logarithms is very small.

For example, an equation resembling

\[ a^m-b^n=1 \]

suggests that

\[ m\log a-n\log b \]

is small.

If transcendence theory provides a lower bound for every nonzero expression

\[ m\log a-n\log b, \]

then the possible values of \(m\) and \(n\) can be bounded.

Thus transcendence theory can turn qualitative irrationality into
quantitative arithmetic information.

%%%%%%%%%%%%%%%%%%%%%%%%%%%%%%%%%%%%%%%%%%%%%%%%%%%%%%%%%%%%
\subsection{Baker's Theorem}
\label{subsec:baker-theorem}
%%%%%%%%%%%%%%%%%%%%%%%%%%%%%%%%%%%%%%%%%%%%%%%%%%%%%%%%%%%%

Alan Baker developed powerful results concerning linear forms in logarithms
of algebraic numbers.

A simplified conceptual version is:

\begin{theorem}[Baker]
Under suitable algebraic hypotheses, a nonzero linear form in logarithms of
algebraic numbers is transcendental and admits explicit nonzero lower bounds.
\end{theorem}

The precise theory requires careful choices of logarithm branches and
algebraic coefficients.

The major practical consequence is that expressions such as

\[ b_1\log\alpha_1+\cdots+b_n\log\alpha_n \]

cannot be nonzero yet arbitrarily small without arithmetic restrictions.

%%%%%%%%%%%%%%%%%%%%%%%%%%%%%%%%%%%%%%%%%%%%%%%%%%%%%%%%%%%%
\subsection{Effective Diophantine Bounds}
\label{subsec:effective-diophantine-bounds}
%%%%%%%%%%%%%%%%%%%%%%%%%%%%%%%%%%%%%%%%%%%%%%%%%%%%%%%%%%%%

A theorem is called \emph{effective} when it provides, at least in principle,
computable bounds.

Baker's theory is especially important because it can produce explicit
bounds for solutions of Diophantine equations.

Schematically,

\[
\boxed{
\text{Diophantine equation}
\longrightarrow
\text{linear form in logarithms}
\longrightarrow
\text{lower bound}
\longrightarrow
\text{bound on integer solutions}.
}
\]

This makes transcendence theory computationally useful.

%%%%%%%%%%%%%%%%%%%%%%%%%%%%%%%%%%%%%%%%%%%%%%%%%%%%%%%%%%%%
\subsection{Exponential Diophantine Equations}
\label{subsec:exponential-diophantine-equations}
%%%%%%%%%%%%%%%%%%%%%%%%%%%%%%%%%%%%%%%%%%%%%%%%%%%%%%%%%%%%

An exponential Diophantine equation contains unknown integers in exponents.

Examples include

\[ 2^m-3^n=1 \]

and

\[ x^m-y^n=1. \]

Such equations can be extraordinarily difficult.

Taking logarithms transforms multiplicative growth into additive
expressions:

\[ m\log2\approx n\log3. \]

The problem then becomes one of controlling rational or algebraic
approximations between logarithms.

%%%%%%%%%%%%%%%%%%%%%%%%%%%%%%%%%%%%%%%%%%%%%%%%%%%%%%%%%%%%
\subsection{Connection with Diophantine Approximation}
\label{subsec:transcendence-diophantine-approximation}
%%%%%%%%%%%%%%%%%%%%%%%%%%%%%%%%%%%%%%%%%%%%%%%%%%%%%%%%%%%%

Liouville's theorem already demonstrates the central connection:

\[
\text{algebraic structure}
\Longrightarrow
\text{limits on approximation}.
\]

Conversely,

\[
\text{approximation that is too strong}
\Longrightarrow
\text{transcendence}.
\]

Modern transcendence theory greatly extends this philosophy beyond rational
approximation to include exponential functions, logarithms, and algebraic
relations among several numbers.

%%%%%%%%%%%%%%%%%%%%%%%%%%%%%%%%%%%%%%%%%%%%%%%%%%%%%%%%%%%%
\subsection{Measures of Transcendence}
\label{subsec:transcendence-measures}
%%%%%%%%%%%%%%%%%%%%%%%%%%%%%%%%%%%%%%%%%%%%%%%%%%%%%%%%%%%%

Instead of merely proving that \(\alpha\) is transcendental, one may ask how
closely \(\alpha\) can satisfy polynomial equations with integer
coefficients.

For a polynomial

\[ P(x)\in\Z[x], \]

one studies how small

\[ |P(\alpha)| \]

can become relative to:

\begin{itemize}
    \item the degree of \(P\);
    \item the size of its coefficients.
\end{itemize}

Quantitative lower bounds are called \emph{transcendence measures}.

These generalize irrationality measures.

%%%%%%%%%%%%%%%%%%%%%%%%%%%%%%%%%%%%%%%%%%%%%%%%%%%%%%%%%%%%
\subsection{Irrationality Versus Transcendence Measures}
\label{subsec:irrationality-vs-transcendence-measures}
%%%%%%%%%%%%%%%%%%%%%%%%%%%%%%%%%%%%%%%%%%%%%%%%%%%%%%%%%%%%

An irrationality measure controls approximations

\[ \left|\alpha-\frac pq\right|. \]

A transcendence measure controls more general expressions

\[ |P(\alpha)|. \]

Thus:

\[
\boxed{
\text{rational approximation}
\subset
\text{polynomial approximation}.
}
\]

Transcendence measures ask whether \(\alpha\) can behave approximately like
a root of a low-complexity integer polynomial.

%%%%%%%%%%%%%%%%%%%%%%%%%%%%%%%%%%%%%%%%%%%%%%%%%%%%%%%%%%%%
\subsection{Normal Numbers}
\label{subsec:normal-numbers-transcendence}
%%%%%%%%%%%%%%%%%%%%%%%%%%%%%%%%%%%%%%%%%%%%%%%%%%%%%%%%%%%%

Another property sometimes associated with complicated real numbers is
normality.

A real number is normal in base \(b\) if every finite block of base-\(b\)
digits appears with the expected limiting frequency.

Normality and transcendence are different properties.

A number may be known to be transcendental without being known to be normal.

Thus:

\[
\boxed{
\text{transcendence is algebraic;}
\qquad
\text{normality is statistical/digital}.
}
\]

%%%%%%%%%%%%%%%%%%%%%%%%%%%%%%%%%%%%%%%%%%%%%%%%%%%%%%%%%%%%
\subsection{Almost All Numbers}
\label{subsec:almost-all-transcendental}
%%%%%%%%%%%%%%%%%%%%%%%%%%%%%%%%%%%%%%%%%%%%%%%%%%%%%%%%%%%%

Because the algebraic numbers are countable, they have Lebesgue measure
zero.

Consequently, almost every real number is transcendental.

In measure-theoretic language,

\[ \boxed{\lambda(\overline{\Q}\cap\R)=0.} \]

Here \(\lambda\) denotes Lebesgue measure, introduced in
\cref{sec:measure-theory}.

This result is striking:

\begin{center}
Transcendental numbers are overwhelmingly common, yet explicit natural
examples can be difficult to analyze.
\end{center}

%%%%%%%%%%%%%%%%%%%%%%%%%%%%%%%%%%%%%%%%%%%%%%%%%%%%%%%%%%%%
\subsection{Algebraic Independence of \(e\) and \(\pi\)}
\label{subsec:e-pi-algebraic-independence}
%%%%%%%%%%%%%%%%%%%%%%%%%%%%%%%%%%%%%%%%%%%%%%%%%%%%%%%%%%%%

We know individually that

\[ e \]

and

\[ \pi \]

are transcendental.

A much stronger question asks whether they are algebraically independent.

That would mean there is no nonzero polynomial

\[ P(x,y)\in\Q[x,y] \]

such that

\[ P(e,\pi)=0. \]

This is not presently known.

Consequently, seemingly simple expressions involving both \(e\) and \(\pi\)
can remain mysterious.

%%%%%%%%%%%%%%%%%%%%%%%%%%%%%%%%%%%%%%%%%%%%%%%%%%%%%%%%%%%%
\subsection{Known and Unknown Constants}
\label{subsec:known-unknown-transcendence}
%%%%%%%%%%%%%%%%%%%%%%%%%%%%%%%%%%%%%%%%%%%%%%%%%%%%%%%%%%%%

It is useful to separate established results from open problems.

\[
\begin{array}{c|c}
\textbf{Number} & \textbf{Known status}\\
\hline
e & \text{transcendental}\\
\pi & \text{transcendental}\\
e^\pi & \text{transcendental}\\
2^{\sqrt2} & \text{transcendental}\\
\log2 & \text{transcendental}\\
e+\pi & \text{not known to be rational, irrational, algebraic, or transcendental}\\
e\pi & \text{not known to be rational, irrational, algebraic, or transcendental}\\
\pi^e & \text{transcendence not known}
\end{array}
\]

\begin{warningbox}
The fact that two constants are individually transcendental does not
determine the arithmetic nature of their sum, product, or arbitrary powers.
\end{warningbox}

%%%%%%%%%%%%%%%%%%%%%%%%%%%%%%%%%%%%%%%%%%%%%%%%%%%%%%%%%%%%
\subsection{Schanuel's Conjecture}
\label{subsec:schanuel-conjecture}
%%%%%%%%%%%%%%%%%%%%%%%%%%%%%%%%%%%%%%%%%%%%%%%%%%%%%%%%%%%%

Many unresolved transcendence questions would follow from a major conjecture
concerning exponentials.

\begin{conjecture}[Schanuel]
Let

\[ z_1,\ldots,z_n\in\C \]

be linearly independent over \(\Q\).

Then the field

\[ \Q(z_1,\ldots,z_n,e^{z_1},\ldots,e^{z_n}) \]

has transcendence degree at least \(n\) over \(\Q\).
\end{conjecture}

This conjecture predicts extensive algebraic independence between numbers
and their exponentials.

It remains unproved.

%%%%%%%%%%%%%%%%%%%%%%%%%%%%%%%%%%%%%%%%%%%%%%%%%%%%%%%%%%%%
\subsection{Interpreting Schanuel's Conjecture}
\label{subsec:interpreting-schanuel}
%%%%%%%%%%%%%%%%%%%%%%%%%%%%%%%%%%%%%%%%%%%%%%%%%%%%%%%%%%%%

The conjecture says, roughly, that linearly independent inputs to the
exponential function should create many algebraically independent numbers.

Schematically,

\[
\boxed{
\text{linear independence of exponents}
\Longrightarrow
\text{large algebraic independence of exponentials}.
}
\]

It would unify and greatly extend several classical transcendence theorems.

%%%%%%%%%%%%%%%%%%%%%%%%%%%%%%%%%%%%%%%%%%%%%%%%%%%%%%%%%%%%
\subsection{The Exponential Function as a Bridge}
\label{subsec:exponential-function-bridge}
%%%%%%%%%%%%%%%%%%%%%%%%%%%%%%%%%%%%%%%%%%%%%%%%%%%%%%%%%%%%

The exponential function repeatedly connects algebraic and transcendental
structures.

Examples include

\[
e^\alpha,
\qquad
\alpha^\beta=e^{\beta\log\alpha},
\qquad
e^{i\pi}=-1.
\]

Thus exponential identities can transform questions about:

\[
\text{powers}
\longleftrightarrow
\text{logarithms}
\longleftrightarrow
\text{algebraic relations}.
\]

This is one reason complex analysis plays such an important role in
transcendence theory.

%%%%%%%%%%%%%%%%%%%%%%%%%%%%%%%%%%%%%%%%%%%%%%%%%%%%%%%%%%%%
\subsection{A Conceptual Map of Classical Results}
\label{subsec:transcendence-classical-map}
%%%%%%%%%%%%%%%%%%%%%%%%%%%%%%%%%%%%%%%%%%%%%%%%%%%%%%%%%%%%

\begin{centereddiagram}
\begin{tikzpicture}[
node distance=1.25cm and 1.6cm,
every node/.style={draw,rounded corners,align=center,inner sep=5pt}
]
\node (L) {Liouville\\approximation};
\node[right=of L] (LC) {Liouville\\numbers};
\node[right=of LC] (T) {explicit\\transcendentals};

\node[below=of L] (HL) {Hermite--Lindemann};
\node[right=of HL] (E) {\(e\)};
\node[right=of E] (PI) {\(\pi\)};

\node[below=of HL] (GS) {Gelfond--Schneider};
\node[right=of GS] (Powers) {algebraic\\irrational powers};

\node[below=of GS] (B) {Baker theory};
\node[right=of B] (Logs) {linear forms\\in logarithms};
\node[right=of Logs] (D) {effective\\Diophantine bounds};

\draw[-{Latex[length=2.2mm]}] (L)--(LC);
\draw[-{Latex[length=2.2mm]}] (LC)--(T);
\draw[-{Latex[length=2.2mm]}] (HL)--(E);
\draw[-{Latex[length=2.2mm]}] (E)--(PI);
\draw[-{Latex[length=2.2mm]}] (GS)--(Powers);
\draw[-{Latex[length=2.2mm]}] (B)--(Logs);
\draw[-{Latex[length=2.2mm]}] (Logs)--(D);
\end{tikzpicture}
\end{centereddiagram}

The methods differ, but all seek arithmetic restrictions strong enough to
rule out algebraicity.

%%%%%%%%%%%%%%%%%%%%%%%%%%%%%%%%%%%%%%%%%%%%%%%%%%%%%%%%%%%%
\subsection{Worked Example: A Transcendental Exponential}
\label{subsec:worked-transcendental-exponential}
%%%%%%%%%%%%%%%%%%%%%%%%%%%%%%%%%%%%%%%%%%%%%%%%%%%%%%%%%%%%

\begin{workedexamplebox}{Classify \(e^{\sqrt3}\)}

Determine whether

\[ e^{\sqrt3} \]

is algebraic or transcendental.

\begin{tabularx}{\textwidth}{@{}>{\raggedright\arraybackslash}p{0.43\textwidth}X@{}}
Classify the exponent. & \(\sqrt3\) is nonzero and algebraic.\\
Apply Hermite--Lindemann. & \(e^\alpha\) is transcendental for nonzero algebraic \(\alpha\).\\
Substitute \(\alpha=\sqrt3\). & \(e^{\sqrt3}\) is transcendental.
\end{tabularx}

\begin{answerbox}
\[ \boxed{e^{\sqrt3}\text{ is transcendental}.} \]
\end{answerbox}
\end{workedexamplebox}

%%%%%%%%%%%%%%%%%%%%%%%%%%%%%%%%%%%%%%%%%%%%%%%%%%%%%%%%%%%%
\subsection{Worked Example: A Logarithm}
\label{subsec:worked-transcendental-logarithm}
%%%%%%%%%%%%%%%%%%%%%%%%%%%%%%%%%%%%%%%%%%%%%%%%%%%%%%%%%%%%

\begin{workedexamplebox}{Classify \(\log 5\)}

Suppose the real logarithm

\[ \log5 \]

were algebraic.

\begin{tabularx}{\textwidth}{@{}>{\raggedright\arraybackslash}p{0.43\textwidth}X@{}}
Assume \(\log5\) is algebraic. & \(\log5\neq0\)\\
Apply Hermite--Lindemann. & \(e^{\log5}\) would be transcendental.\\
Evaluate the exponential. & \(e^{\log5}=5\)\\
But \(5\) is algebraic. & Contradiction
\end{tabularx}

\begin{answerbox}
\[ \boxed{\log5\text{ is transcendental}.} \]
\end{answerbox}
\end{workedexamplebox}

%%%%%%%%%%%%%%%%%%%%%%%%%%%%%%%%%%%%%%%%%%%%%%%%%%%%%%%%%%%%
\subsection{Worked Example: What a Theorem Cannot Prove}
\label{subsec:worked-transcendence-theorem-limit}
%%%%%%%%%%%%%%%%%%%%%%%%%%%%%%%%%%%%%%%%%%%%%%%%%%%%%%%%%%%%

\begin{workedexamplebox}{Can Gelfond--Schneider Prove \(\pi^e\) Is Transcendental?}

\begin{tabularx}{\textwidth}{@{}>{\raggedright\arraybackslash}p{0.43\textwidth}X@{}}
Check the required base. & It must be algebraic.\\
Actual base. & \(\pi\) is transcendental.\\
Check the required exponent. & It must be algebraic irrational.\\
Actual exponent. & \(e\) is transcendental.\\
Conclusion. & The theorem does not apply.
\end{tabularx}

\begin{answerbox}
\[ \boxed{\text{Gelfond--Schneider does not determine }\pi^e.} \]
\end{answerbox}
\end{workedexamplebox}

%%%%%%%%%%%%%%%%%%%%%%%%%%%%%%%%%%%%%%%%%%%%%%%%%%%%%%%%%%%%
\subsection{Common Errors}
\label{subsec:transcendental-number-theory-common-errors}
%%%%%%%%%%%%%%%%%%%%%%%%%%%%%%%%%%%%%%%%%%%%%%%%%%%%%%%%%%%%

\subsubsection{Confusing Irrational and Transcendental}

Every transcendental real number is irrational.

But not every irrational number is transcendental.

For example,

\[ \sqrt2 \]

is irrational but algebraic.

%%%%%%%%%%%%%%%%%%%%%%%%%%%%%%%%%%%%%%%%%%%%%%%%%%%%%%%%%%%%
\subsubsection{Assuming Complicated Decimals Are Transcendental}

A nonrepeating decimal is irrational, but irrationality alone does not imply
transcendence.

%%%%%%%%%%%%%%%%%%%%%%%%%%%%%%%%%%%%%%%%%%%%%%%%%%%%%%%%%%%%
\subsubsection{Assuming Every Transcendental Number Is Liouville}

Liouville numbers form only a special class of transcendental numbers.

%%%%%%%%%%%%%%%%%%%%%%%%%%%%%%%%%%%%%%%%%%%%%%%%%%%%%%%%%%%%
\subsubsection{Adding Transcendental Numbers Carelessly}

From

\[ \alpha,\beta\text{ transcendental} \]

one cannot conclude

\[ \alpha+\beta\text{ transcendental}. \]

%%%%%%%%%%%%%%%%%%%%%%%%%%%%%%%%%%%%%%%%%%%%%%%%%%%%%%%%%%%%
\subsubsection{Multiplying Transcendental Numbers Carelessly}

Likewise,

\[ \alpha\beta \]

may be algebraic even if both factors are transcendental.

%%%%%%%%%%%%%%%%%%%%%%%%%%%%%%%%%%%%%%%%%%%%%%%%%%%%%%%%%%%%
\subsubsection{Misapplying Gelfond--Schneider}

The base must be algebraic and different from \(0\) and \(1\).

The exponent must be algebraic and irrational.

%%%%%%%%%%%%%%%%%%%%%%%%%%%%%%%%%%%%%%%%%%%%%%%%%%%%%%%%%%%%
\subsubsection{Misapplying Hermite--Lindemann}

The theorem requires a nonzero algebraic exponent.

It does not say that \(e^\alpha\) is transcendental for every transcendental
\(\alpha\).

Indeed,

\[ e^{\log2}=2. \]

%%%%%%%%%%%%%%%%%%%%%%%%%%%%%%%%%%%%%%%%%%%%%%%%%%%%%%%%%%%%
\subsubsection{Confusing Linear and Algebraic Independence}

Linear independence rules out relations of the form

\[ a_1x_1+\cdots+a_nx_n=0. \]

Algebraic independence rules out every nonzero polynomial relation.

Algebraic independence is much stronger.

%%%%%%%%%%%%%%%%%%%%%%%%%%%%%%%%%%%%%%%%%%%%%%%%%%%%%%%%%%%%
\subsubsection{Treating Open Problems as Established Results}

Expressions involving \(e\) and \(\pi\) can remain unresolved even though
both constants are individually transcendental.

%%%%%%%%%%%%%%%%%%%%%%%%%%%%%%%%%%%%%%%%%%%%%%%%%%%%%%%%%%%%
\subsection{Applications and Connections}
\label{subsec:transcendental-number-theory-applications}
%%%%%%%%%%%%%%%%%%%%%%%%%%%%%%%%%%%%%%%%%%%%%%%%%%%%%%%%%%%%

\begin{applicationbox}
\textbf{Diophantine Approximation.}
Liouville's theorem uses rational approximation to distinguish algebraic
numbers from certain transcendental numbers.

\medskip

\textbf{Algebraic Number Theory.}
Transcendence theory begins with the arithmetic structure of algebraic
numbers, number fields, minimal polynomials, and logarithms of algebraic
numbers.

\medskip

\textbf{Complex Analysis.}
The exponential function, logarithms, analytic functions, and complex
identities such as \(e^{i\pi}=-1\) are central to classical transcendence
proofs.

\medskip

\textbf{Diophantine Equations.}
Baker's theory converts lower bounds for linear forms in logarithms into
effective bounds for integer solutions.

\medskip

\textbf{Geometry.}
The transcendence of \(\pi\) proves the impossibility of squaring the circle
with straightedge and compass.

\medskip

\textbf{Field Theory.}
Algebraic independence and transcendence degree measure how far field
extensions extend beyond algebraic equations.

\medskip

\textbf{Arithmetic Geometry.}
Periods, logarithms, elliptic functions, and special values lead to deeper
transcendence and algebraic-independence questions.

\medskip

\textbf{Mathematical Logic.}
Cardinality arguments establish the abundance of transcendental numbers even
without explicitly constructing them.

\medskip

\textbf{Computational Number Theory.}
Effective transcendence estimates can bound searches for solutions of
exponential Diophantine equations.
\end{applicationbox}

%%%%%%%%%%%%%%%%%%%%%%%%%%%%%%%%%%%%%%%%%%%%%%%%%%%%%%%%%%%%
\subsection{Exercises}
\label{subsec:transcendental-number-theory-exercises}
%%%%%%%%%%%%%%%%%%%%%%%%%%%%%%%%%%%%%%%%%%%%%%%%%%%%%%%%%%%%

\begin{exercise}[1]
Classify each number as rational, algebraic irrational, or known
transcendental:

\[ \frac57,\qquad\sqrt3,\qquad e,\qquad\pi. \]
\end{exercise}

\begin{exercise}[1]
Give an integer polynomial having

\[ \sqrt7 \]

as a root.
\end{exercise}

\begin{exercise}[1]
Give an integer polynomial having

\[ \frac{2+\sqrt3}{5} \]

as a root.
\end{exercise}

\begin{exercise}[2]
Explain why every rational number is algebraic.
\end{exercise}

\begin{exercise}[2]
Explain why the algebraic numbers form a countable set.
\end{exercise}

\begin{exercise}[2]
Use cardinality to explain why transcendental real numbers exist.
\end{exercise}

\begin{exercise}[2]
Why does the cardinality argument fail to prove directly that \(\pi\) is
transcendental?
\end{exercise}

\begin{exercise}[2]
State Liouville's approximation theorem.
\end{exercise}

\begin{exercise}[2]
Explain why every Liouville number is transcendental.
\end{exercise}

\begin{exercise}[2]
Write the first four terms of

\[ \sum_{n=1}^{\infty}10^{-n!}. \]
\end{exercise}

\begin{exercise}[3]
Explain why truncations of Liouville's constant give unusually strong
rational approximations.
\end{exercise}

\begin{exercise}[2]
Explain why Liouville's theorem does not prove that every transcendental
number is a Liouville number.
\end{exercise}

\begin{exercise}[1]
Use Hermite--Lindemann to classify

\[ e^2. \]
\end{exercise}

\begin{exercise}[1]
Use Hermite--Lindemann to classify

\[ e^{\sqrt5}. \]
\end{exercise}

\begin{exercise}[2]
Use Hermite--Lindemann to prove that

\[ \log2 \]

is transcendental.
\end{exercise}

\begin{exercise}[2]
Explain how Euler's identity and Hermite--Lindemann imply the transcendence
of \(\pi\).
\end{exercise}

\begin{exercise}[2]
Explain why the transcendence of \(\pi\) implies that squaring the circle is
impossible using straightedge and compass.
\end{exercise}

\begin{exercise}[2]
State the Lindemann--Weierstrass theorem conceptually.
\end{exercise}

\begin{exercise}[2]
Explain the difference between linear independence and algebraic
independence.
\end{exercise}

\begin{exercise}[2]
Show that if \(\alpha\) is transcendental, then

\[ \alpha,\alpha^2 \]

are algebraically dependent.
\end{exercise}

\begin{exercise}[1]
Use Gelfond--Schneider to prove that

\[ 2^{\sqrt3} \]

is transcendental.
\end{exercise}

\begin{exercise}[1]
Use Gelfond--Schneider to classify

\[ 7^{\sqrt2}. \]
\end{exercise}

\begin{exercise}[2]
Explain why Gelfond--Schneider does not directly apply to

\[ \pi^{\sqrt2}. \]
\end{exercise}

\begin{exercise}[2]
Explain why Gelfond--Schneider does not directly apply to

\[ 2^\pi. \]
\end{exercise}

\begin{exercise}[2]
Explain why Gelfond--Schneider does not determine

\[ \pi^e. \]
\end{exercise}

\begin{exercise}[2]
Prove that if \(\alpha\neq0\) is algebraic and \(\beta\) is transcendental,
then

\[ \alpha\beta \]

is transcendental.
\end{exercise}

\begin{exercise}[2]
Prove that if \(\alpha\) is algebraic and \(\beta\) is transcendental, then

\[ \alpha+\beta \]

is transcendental.
\end{exercise}

\begin{exercise}[1]
Determine whether the following conclusion is valid:

\[
e,\pi\text{ transcendental}
\quad\Longrightarrow\quad
e+\pi\text{ transcendental}.
\]

Explain.
\end{exercise}

\begin{exercise}[1]
Determine whether the following conclusion is valid:

\[
e,\pi\text{ transcendental}
\quad\Longrightarrow\quad
e\pi\text{ transcendental}.
\]

Explain.
\end{exercise}

\begin{exercise}[2]
Give two transcendental numbers whose sum is algebraic.
\end{exercise}

\begin{exercise}[2]
Give two transcendental numbers whose product is algebraic.
\end{exercise}

\begin{exercise}[2]
Define algebraic independence over \(\Q\).
\end{exercise}

\begin{exercise}[2]
Explain why

\[ t\qquad\text{and}\qquad t^3 \]

are algebraically dependent whenever \(t\) is transcendental.
\end{exercise}

\begin{exercise}[2]
What is meant by a linear form in logarithms?
\end{exercise}

\begin{exercise}[2]
Identify the coefficients and logarithms in

\[ \Lambda=7\log2-4\log3. \]
\end{exercise}

\begin{exercise}[2]
Explain why an equation such as

\[ 2^m\approx3^n \]

naturally leads to the expression

\[ m\log2-n\log3. \]
\end{exercise}

\begin{exercise}[3]
Explain conceptually how a lower bound for

\[ |m\log2-n\log3| \]

could restrict possible integer solutions to an exponential Diophantine
equation.
\end{exercise}

\begin{exercise}[2]
State Baker's theorem conceptually.
\end{exercise}

\begin{exercise}[2]
Explain what is meant by an effective result.
\end{exercise}

\begin{exercise}[2]
Distinguish an irrationality measure from a transcendence measure.
\end{exercise}

\begin{exercise}[2]
Explain why almost every real number is transcendental with respect to
Lebesgue measure.
\end{exercise}

\begin{exercise}[2]
Explain why normality and transcendence are logically different properties.
\end{exercise}

\begin{exercise}[2]
Which of the following are known to be transcendental?

\[
e,\qquad
\pi,\qquad
e^\pi,\qquad
2^{\sqrt2},\qquad
e+\pi.
\]
\end{exercise}

\begin{exercise}[2]
Which theorem proves the transcendence of

\[ 5^{\sqrt7}? \]
\end{exercise}

\begin{exercise}[2]
Which theorem implies the transcendence of

\[ e^{\sqrt7}? \]
\end{exercise}

\begin{exercise}[3]
Explain the conceptual chain

\[
\text{exceptionally strong rational approximation}
\longrightarrow
\text{Liouville's theorem}
\longrightarrow
\text{transcendence}.
\]
\end{exercise}

\begin{exercise}[3]
Explain the conceptual chain

\[
e^{i\pi}=-1
\longrightarrow
\text{Hermite--Lindemann}
\longrightarrow
\pi\text{ transcendental}.
\]
\end{exercise}

\begin{exercise}[3]
Explain the conceptual chain

\[
\text{algebraic base}
+
\text{algebraic irrational exponent}
\longrightarrow
\text{Gelfond--Schneider}
\longrightarrow
\text{transcendental power}.
\]
\end{exercise}

\begin{exercise}[3]
Explain the conceptual chain

\[
\text{exponential Diophantine equation}
\longrightarrow
\text{linear form in logarithms}
\longrightarrow
\text{Baker-type lower bound}
\longrightarrow
\text{finite search}.
\]
\end{exercise}

\begin{exercise}[3]
Explain why proving that \(e\) and \(\pi\) are each transcendental is weaker
than proving that they are algebraically independent.
\end{exercise}

\begin{exercise}[3]
State Schanuel's conjecture conceptually and explain why it would represent a
major strengthening of classical transcendence theory.
\end{exercise}

%%%%%%%%%%%%%%%%%%%%%%%%%%%%%%%%%%%%%%%%%%%%%%%%%%%%%%%%%%%%
\subsection{Conceptual Summary}
\label{subsec:transcendental-number-theory-summary}
%%%%%%%%%%%%%%%%%%%%%%%%%%%%%%%%%%%%%%%%%%%%%%%%%%%%%%%%%%%%

A number is algebraic when it satisfies a nonzero polynomial equation with
rational coefficients.

A number that satisfies no such equation is transcendental.

Thus

\[
\boxed{
\C=
\{\text{algebraic numbers}\}
\cup
\{\text{transcendental numbers}\}.
}
\]

The algebraic numbers are countable, while the real numbers are uncountable.

Therefore transcendental numbers are abundant.

However, proving that a specific naturally occurring number is
transcendental can be difficult.

Liouville's theorem provides one of the first explicit methods:

\[
\boxed{
\text{too well approximated by rationals}
\Longrightarrow
\text{not algebraic}.
}
\]

This produces Liouville numbers, including

\[ L=\sum_{n=1}^{\infty}10^{-n!}. \]

The Hermite--Lindemann theorem states that

\[ \boxed{\alpha\neq0\text{ algebraic}\Longrightarrow e^\alpha
\text{ transcendental}.} \]

It implies the transcendence of \(e\).

Combined with

\[ e^{i\pi}=-1, \]

it also yields the transcendence of \(\pi\).

The Lindemann--Weierstrass theorem greatly strengthens these exponential
transcendence results.

The Gelfond--Schneider theorem states that, under its standard hypotheses,

\[
\boxed{
\alpha\text{ algebraic},
\quad
\beta\text{ algebraic irrational}
\Longrightarrow
\alpha^\beta\text{ transcendental}.
}
\]

Consequently,

\[ 2^{\sqrt2} \]

and

\[ e^\pi \]

are transcendental.

Baker's theory studies linear forms

\[ \Lambda=b_1\log\alpha_1+\cdots+b_n\log\alpha_n \]

and provides powerful lower bounds for nonzero forms.

These bounds transform transcendence theory into a tool for solving
Diophantine equations.

Transcendence of individual numbers must be distinguished from algebraic
independence.

Even though \(e\) and \(\pi\) are individually transcendental, many simple
questions involving both constants remain unresolved.

The central principle is

\[
\boxed{
\text{transcendence theory measures the limits of algebraic relations among
important numbers}.
}
\]

%%%%%%%%%%%%%%%%%%%%%%%%%%%%%%%%%%%%%%%%%%%%%%%%%%%%%%%%%%%%
\subsection{Section Connections}
\label{subsec:transcendental-number-theory-connections}
%%%%%%%%%%%%%%%%%%%%%%%%%%%%%%%%%%%%%%%%%%%%%%%%%%%%%%%%%%%%

\chapterconnection{
Transcendental Number Theory studies numbers that cannot satisfy polynomial
equations with rational coefficients. Diophantine Approximation supplies the
first major transcendence mechanism through Liouville's theorem. Algebraic
Number Theory provides the theory of algebraic numbers, number fields, and
logarithms that form the arithmetic input to later transcendence theorems.
Complex Analysis supplies exponential and logarithmic functions, allowing
Hermite--Lindemann and Lindemann--Weierstrass to connect algebraic inputs
with transcendental outputs. Gelfond--Schneider resolves transcendence
questions involving algebraic irrational powers, while Baker's theory of
linear forms in logarithms provides effective methods for exponential
Diophantine equations. Algebraic independence and transcendence degree
connect the subject with Field Theory and Algebraic Geometry. These themes
lead naturally into Arithmetic Geometry, where algebraic varieties,
Diophantine equations, heights, local methods, and analytic functions are
combined into a unified arithmetic framework.
}

%%%%%%%%%%%%%%%%%%%%%%%%%%%%%%%%%%%%%%%%%%%%%%%%%%%%%%%%%%%%
% SECTION SOURCE TRACKING
%%%%%%%%%%%%%%%%%%%%%%%%%%%%%%%%%%%%%%%%%%%%%%%%%%%%%%%%%%%%

%------------------------------------------------------------
% 2.8 Transcendental Number Theory
%------------------------------------------------------------
%
% Source basis:
%   - Standard classical transcendental number theory.
%
% Used for:
%   - algebraic and transcendental numbers;
%   - countability of algebraic numbers;
%   - existence of transcendental numbers;
%   - Liouville's approximation theorem;
%   - Liouville numbers and Liouville's constant;
%   - Hermite's theorem on e;
%   - Hermite--Lindemann theorem;
%   - transcendence of pi;
%   - squaring the circle;
%   - Lindemann--Weierstrass theorem;
%   - linear independence of exponentials;
%   - algebraic independence;
%   - transcendence degree;
%   - Gelfond--Schneider theorem;
%   - Hilbert's seventh problem;
%   - transcendence of 2^{sqrt(2)};
%   - transcendence of e^pi;
%   - logarithms of algebraic numbers;
%   - linear forms in logarithms;
%   - Baker's theorem overview;
%   - effective Diophantine bounds;
%   - transcendence measures;
%   - known and open transcendence questions;
%   - Schanuel's conjecture.
%
% Notes:
%   - Liouville approximation is referenced rather than
%     fully redeveloped because its canonical introduction is
%     in Diophantine Approximation.
%   - Algebraic-number structure belongs canonically to
%     Algebraic Number Theory.
%   - Complex exponential and logarithmic analysis belongs
%     canonically to Complex Analysis.
%   - Full proofs of Lindemann--Weierstrass,
%     Gelfond--Schneider, and Baker's theorem require
%     substantially more advanced analytic and algebraic
%     machinery than is developed here.
%   - Algebraic independence is introduced conceptually here
%     and developed further through Field Theory and advanced
%     arithmetic topics.
%   - Schanuel's conjecture is clearly presented as an open
%     conjecture rather than an established theorem.
%
%%%%%%%%%%%%%%%%%%%%%%%%%%%%%%%%%%%%%%%%%%%%%%%%%%%%%%%%%%%%